\begin{lemma}[$p$-harmonicity transfers to the union of an increasing exhaustion]
Let $1<p<2$, let $(\Omega_n)_{n\in\mathbb{N}}$ be an increasing family of subsets of
$\mathrm{Plane}$, that is $\Omega_n\subseteq\Omega_m$ whenever $n\le m$, let
$u\colon\mathrm{Plane}\to\mathbb{R}$ and $g\colon\mathrm{Plane}\to\mathrm{Plane}$, and assume
that for every $n$ the function $u$ is $p$-harmonic in $\Omega_n$ with weak gradient $g$ --- the
same $g$ for every $n$.  Then $u$ is $p$-harmonic in $\Omega_\infty:=\bigcup_{n}\Omega_n$ with
weak gradient $g$.
\end{lemma}
\begin{proof}
We refer to the six conditions in the definition of a $p$-harmonic function with weak gradient,
\noderef{IsPHarm}, by their numbers (1)--(6), and write $A_p$ for the $p$-flux \noderef{PFlux}.
By condition (1) applied to the datum for index $n$, each $\Omega_n$ is open.

\medskip\noindent\emph{Step 1: a compact subset of $\Omega_\infty$ lies in a single
$\Omega_N$.}
Let $K\subseteq\Omega_\infty$ be compact.  The sets $\Omega_n$ are open and cover $K$, so by
compactness there are finitely many indices $n_1,\ldots,n_j$ with
$K\subseteq\Omega_{n_1}\cup\cdots\cup\Omega_{n_j}$.  If $j=0$ then $K=\emptyset\subseteq\Omega_0$
and we may take $N:=0$.  Otherwise put $N:=\max\{n_1,\ldots,n_j\}$.  For each $i$ we have
$n_i\le N$, so $\Omega_{n_i}\subseteq\Omega_N$ by the increasing hypothesis; hence
\begin{equation}\label{eq:pue-one}
  K\subseteq\Omega_{n_1}\cup\cdots\cup\Omega_{n_j}\subseteq\Omega_N .
\end{equation}

\medskip\noindent\emph{Step 2: conditions (1)--(4) for $\Omega_\infty$.}

(1) Each $\Omega_n$ is open, and a union of open sets is open, so $\Omega_\infty$ is open.  In
particular $\Omega_\infty$ is Lebesgue measurable, by clause (2) of the background proposition
on integrability of continuous and compactly supported functions.

(2) Let $x\in\Omega_\infty$, say $x\in\Omega_n$.  By condition (2) for index $n$, the function
$u$ is continuous on $\Omega_n$, and $\Omega_n$ is open, hence a neighbourhood of $x$; so $u$ is
continuous at $x$, by the last clause of the background proposition on uniform limits of
continuous functions.  Thus $u$ is continuous at every point of $\Omega_\infty$, hence
continuous on $\Omega_\infty$.

(3) The index set $\mathbb{N}$ is countable, and by condition (3) for each index $n$ the
function $g$ is almost everywhere strongly measurable with respect to Lebesgue measure
restricted to $\Omega_n$.  Clause (4) of the background proposition on restriction and
congruence --- the countable-union criterion --- therefore gives that $g$ is almost everywhere
strongly measurable with respect to Lebesgue measure restricted to
$\Omega_\infty=\bigcup_n\Omega_n$.

(4) Let $K\subseteq\Omega_\infty$ be compact.  By Step 1 there is $N$ with
$K\subseteq\Omega_N$, and condition (4) for index $N$ gives that
$x\mapsto\lvert g(x)\rvert^{p}$ is integrable on $K$.

\medskip\noindent\emph{Step 3: a localization estimate.}
Let $W\subseteq\mathrm{Plane}$ be open and suppose $u$ is continuous on $W$, $g$ is almost
everywhere strongly measurable for Lebesgue measure restricted to $W$, and
$x\mapsto\lvert g(x)\rvert^{p}$ is integrable on every compact subset of $W$.  Let $\varphi$ be
of class $C^\infty$ with compact support and $\operatorname{supp}\varphi\subseteq W$, and put
$K_0:=\operatorname{supp}\varphi$, a compact subset of $W$, together with
$B:=\sup_y\lvert D\varphi(y)\rvert<\infty$ and $B':=\sup_y\lvert\varphi(y)\rvert<\infty$, finite
because $\varphi$ and $D\varphi$ are continuous and vanish off $K_0$ (on the open set
$\mathrm{Plane}\setminus K_0$ the function $\varphi$ vanishes identically, so all its
derivatives do too).  Then each of
\[
  u\,D\varphi,\qquad \varphi\,g,\qquad \bigl\langle A_p(g),D\varphi\bigr\rangle
\]
vanishes off $K_0$, is integrable on $W$, and has the same integral over $W$ as over $K_0$.

For the first: $W$ is measurable, $u$ is continuous on $W$ and $D\varphi$ on all of
$\mathrm{Plane}$, so $u\,D\varphi$ is continuous on $W$, hence bounded and integrable on the
compact $K_0\subseteq W$ by clause (3) of the integrability proposition; vanishing off $K_0$,
clause (4) of that proposition gives integrability on $W$ and equality of the two integrals.
For the second: taking $a=1$ in clause (4) of the background proposition on elementary
properties of real powers (legitimate as $1\le p$) gives
$\lvert\eta\rvert\le1+\lvert\eta\rvert^{p}$ for every vector $\eta$, hence
$\lvert\varphi\,g\rvert\le B'(1+\lvert g\rvert^{p})$, whose right side is integrable on $K_0$
since $\lvert g\rvert^{p}$ is and $\lvert K_0\rvert<\infty$; $\varphi\,g$ is almost everywhere
strongly measurable on $W$, so clauses (3) and (4) of the integrability proposition finish it.
For the third: by \noderef{PFluxContinuity} the flux $A_p$ is continuous on $\mathrm{Plane}$ with
$\lvert A_p(\eta)\rvert=\lvert\eta\rvert^{\,p-1}$ for every $\eta$, so $A_p\circ g$ is almost
everywhere strongly measurable on $W$ (composition with a continuous map preserves that, by the
background proposition on pointwise limits), and by Cauchy--Schwarz together with clause (4) of
the elementary-powers proposition with $a=p-1\in[0,p]$,
\[
  \bigl\lvert\bigl\langle A_p(g(y)),D\varphi(y)\bigr\rangle\bigr\rvert
  \le B\,\lvert g(y)\rvert^{\,p-1}\le B\bigl(1+\lvert g(y)\rvert^{p}\bigr),
\]
integrable on $K_0$ as before; clauses (3) and (4) again finish it.

\medskip\noindent\emph{Step 4: conditions (5) and (6) for $\Omega_\infty$.}
Let $\varphi$ be of class $C^\infty$ with compact support and
$\operatorname{supp}\varphi\subseteq\Omega_\infty$, and put $K_0:=\operatorname{supp}\varphi$,
which is compact.  By Step 1 applied to $K_0$ there is $N$ with $K_0\subseteq\Omega_N$.  Thus
$\varphi$ is an admissible test function for $\Omega_N$, and conditions (5) and (6) for index
$N$ give
\begin{equation}\label{eq:pue-N}
  \int_{\Omega_N}u\,D\varphi\,dx=-\int_{\Omega_N}\varphi\,g\,dx,
  \qquad
  \int_{\Omega_N}\bigl\langle A_p(g),D\varphi\bigr\rangle\,dx=0 .
\end{equation}
Step 3 applies with $W=\Omega_N$, whose three hypotheses are conditions (2), (3), (4) for index
$N$, and also with $W=\Omega_\infty$, whose three hypotheses are the conditions (2), (3), (4)
established in Step 2.  In both cases $\operatorname{supp}\varphi=K_0$ is contained in $W$.
Hence all six integrals appearing in \eqref{eq:pue-N} and in the corresponding statements over
$\Omega_\infty$ reduce to integrals over $K_0$:
\[
  \int_{\Omega_\infty}u\,D\varphi=\int_{K_0}u\,D\varphi=\int_{\Omega_N}u\,D\varphi,
  \qquad
  \int_{\Omega_\infty}\varphi\,g=\int_{K_0}\varphi\,g=\int_{\Omega_N}\varphi\,g,
\]
\[
  \int_{\Omega_\infty}\bigl\langle A_p(g),D\varphi\bigr\rangle
  =\int_{K_0}\bigl\langle A_p(g),D\varphi\bigr\rangle
  =\int_{\Omega_N}\bigl\langle A_p(g),D\varphi\bigr\rangle .
\]
Substituting into \eqref{eq:pue-N},
\[
  \int_{\Omega_\infty}u\,D\varphi\,dx=-\int_{\Omega_\infty}\varphi\,g\,dx,
  \qquad
  \int_{\Omega_\infty}\bigl\langle A_p(g),D\varphi\bigr\rangle\,dx=0 ,
\]
which are conditions (5) and (6) for $u$, $g$ on $\Omega_\infty$.

\medskip\noindent All six conditions hold, so $u$ is $p$-harmonic in $\Omega_\infty$ with weak
gradient $g$.
\end{proof}
